O mundo, neste momento, está passando pelo que pode ser chamado como a era da digitalização, 
havendo severas transformações no modo de transmissão de dados e informações \cite{Tedjojuwono2024}. 
Atualmente, a humanidade encontra-se sob influência da expansão da denominada indústria 4.0,
que tem como foco a digitalização dos processos de manufatura e produção, mas que engloba 
inclusive o modo de vida e acesso à informações da sociedade sob o braço 
tecnológico denominado \textit{Cidades Inteligentes}.

Neste contexto, a utilização de hardware inteligente deixou de ser mero luxo de grandes
corporações para se tornar o paradigma central na mudança citada anteriormente. Em relação à equipamentos de medição inteligentes,
também classificados como \gls{smart_metering}, hoje já há uma série de avanços que possibilitam o processamento 
de dados de maneira mais veloz e eficiente. Mas essa evolução também acompanhou a necessidade de adaptação de infraestrutura, 
assim agregando altos custos para à transição das tecnologias analógicas para o digital \cite{Stefano2019}.

Como apresentado por \citeonline{Stefano2019}, o \gls{retrofitting} é uma das soluções mais difundidas para contornar o 
alto custo de digitalização de sistemas nativamente analógicos. Tal conceito, tem como objetivo realizar a digitalização 
de sistemas analógicos reaproveitando o hardware já existente. Essa concepção, é amplamente abordada na academia, havendo grandes teses,
trabalhos e dissertações sobre o tema.

Diante de tal abordagem para redução de custos, a aplicabilidade da técnica de \gls{retrofitting} enfrenta um impasse.
Atualmente, na academia encontra-se mais difundido trabalhos que abordam o uso de microcontroladores aliados à tecnologia \gls{lora}
utilizando o método de processamento em borda, como explicado por \citeonline{SilvaVilela2021}, a utilização deste método 
traz grandes vantagens em relação à latência e otimização de consumo energético. Em oposição, como \citeonline{Francesco2024} apresenta, 
o método de processamento em cloud destaca-se por proporcionar maior granularidade de dados, simplificação no desenvolvimento 
de firmware e acurácia ligeiramente melhor em relação ao conceito de \gls{edge_computing}.Todavia, permanece a dúvida de qual 
método de processamento utilizar no desenvolvimento de uma solução \gls{iot} baseado na técnica de \gls{retrofitting}. 
A resolução do problema criará um alicerce, auxiliando o desenvolvimento de futuras ferramentas, que por concepção, apresentará a 
mesma natureza de problemática.
